\documentclass[10pt,leqno]{amsart}
\usepackage{graphicx}
\baselineskip=16pt

\usepackage{indentfirst,csquotes}

\topmargin= .5cm
\textheight= 20cm
\textwidth= 32cc
\baselineskip=16pt

\evensidemargin= .9cm
\oddsidemargin= .9cm

\usepackage{amssymb,amsthm,amsmath}
\usepackage{xcolor,paralist,hyperref,titlesec,fancyhdr,etoolbox}
\newtheorem{theorem}{Theorem}[]
\newtheorem{definition}[theorem]{Definition}
\newtheorem{example}[theorem]{Example}
\newtheorem{lemma}[theorem]{Lemma}
\newtheorem{proposition}[theorem]{Proposition}
\newtheorem{corollary}[theorem]{Corollary}
\newtheorem{conjecture}[theorem]{Conjecture}

\titleformat{\section}[display]{\normalfont\huge\bfseries\centering}{\centering\thesection}{10pt}{\Large}
\titlespacing*{\section}{0pt}{0ex}{0ex}

\hypersetup{ colorlinks=true, linkcolor=black, filecolor=black, urlcolor=black }
\def\proof{\noindent {\it Proof. $\, $}}
\def\endproof{\hfill $\Box$ \vskip 5 pt }

\usepackage{lipsum}

% helpers for the disposition ------------------------------------------------
\newcommand{\todo}[1]{{\color{red}[TODO: #1]}}
\newcommand{\ci}{\perp\!\!\!\perp}          % conditional independence symbol
\newcommand{\nci}{\not\!\perp\!\!\!\perp}

\begin{document}
\title{Causal Discovery Between Blocks} %%%%%%%%%%%%
\author[Initial Surname]{C.B.H. Thorjussen \& K. H. Liland}
\date{\today}
\address{Telemarksforsking}
\email{christian.b.thorjussen@tmforsk.com}
\maketitle

\let\thefootnote\relax
\footnotetext{MSC2020: Primary 00A05, Secondary 00A66.} %%%%%%%%%%

\begin{abstract}
\emph{Draft.} Causal discovery methods such as the PC algorithm operate on single
variables, whereas many data sets in chemometrics, sensory science and the
omics fields come as \emph{blocks}: matrices of many variables measured on the
same objects. We define graphical models whose nodes are blocks, propose a test of
conditional independence between blocks based on sequential orthogonal PLS
(SO-PLS) with a resampling-based null distribution, and use it inside the PC
algorithm to recover the causal structure between blocks. In simulations the test
\todo{state final type I error / power numbers} and the resulting block-level PC
algorithm \todo{state final SHD / recall numbers}. We illustrate the method on
\todo{real data sets}.
\end{abstract} %%%%%%%%%

\bigskip

\noindent

$\,$

$\,$
% =============================================================================
\section{Introduction}
\begin{itemize}
\item \textbf{Motivation.} Multiblock data: several matrices on the same objects
  (spectra, sensory profiles, process settings, consumer data, omics). Questions
  are causal (what drives what), not only predictive.
\item \textbf{Background.} Constraint-based causal discovery (PC, FCI; Spirtes and
  Glymour 1991; Spirtes, Glymour and Scheines 2000) and its reliance on
  conditional independence (CI) tests of single variables. Existing multivariate CI
  tests (kernel, distance-based, CCI) and their limits (Shah and Peters 2020).
\item \textbf{Gap.} (i) No graphical model with blocks as nodes; (ii) no CI test
  between blocks that is fast, handles collinear high-dimensional blocks, and
  respects multiblock structure; (iii) no PC-type algorithm built on such a test.
\item \textbf{Contributions.} (1) Block graphical models (Section 2).
  (2) SO-PLS-based CI test, with fast and residual-based variants (Section 3).
  (3) Study of the effect of block order (Section 5).
  (4) Block-level PC algorithm (Section 6).
  (5) Simulation and real-data evidence (Sections 4, 7, 8).
  (6) Open-source R package (\texttt{CausalBlocks}/\texttt{BlockCI}).
\item \textbf{Two settings.} (A) same samples, different variables (the focus);
  (B) same variables, different samples/environments (invariant prediction;
  brief remarks only).
\item \textbf{Outline of the paper.}
\end{itemize}

% =============================================================================
\section{Block Graphical Modelling}
\begin{itemize}
\item \textbf{Notation.} Block $X_k\in\mathbb{R}^{n\times p_k}$, rows are objects.
  A partition of the variable set into blocks $\mathcal{B}=\{B_1,\dots,B_K\}$.
\item \textbf{Definition (block CI).} $B_i \ci B_j \mid B_S$ means the random
  vectors are conditionally independent. Note this is a statement about whole
  blocks, not about individual columns.
\item \textbf{Block graph.} Nodes are blocks, edges between blocks. Relation to an
  underlying variable-level DAG $\mathcal{G}$: block graph as a \emph{contraction} of
  $\mathcal{G}$ along the partition.
  \begin{itemize}
  \item When is the contracted graph acyclic? Condition on the partition
    (e.g.\ a topological ordering consistent with blocks).
  \item Edges inside a block are ignored (``within-block structure is left
    unmodelled'').
  \end{itemize}
\item \textbf{Markov property and faithfulness at block level.} Block
  d-separation, global Markov property, which block-level independences are implied
  by the variable-level graph. \todo{state and prove: block d-separation $\Leftrightarrow$
  d-separation of all variables in the blocks, i.e.\ the ``composition'' property;
  note it requires faithfulness of the joint distribution, and that composition
  can fail without it.}
\item \textbf{Markov equivalence of block graphs.} v-structures at block level,
  CPDAG of blocks.
\item \textbf{Examples} used throughout: common cause ($Z_1,Z_2\to X,Y$), chain
  ($X\leftarrow Z_1\to Z_2\to Y$), collider, and a four-block graph.
\item \textbf{Assumptions and limits.} Causal sufficiency at block level (latent
  variables inside blocks are fine, shared hidden causes between blocks are not);
  linearity assumptions where the test needs them.
\end{itemize}

% =============================================================================
\section{Conditional Independence Test between Blocks of Data}
\begin{itemize}
\item \textbf{Problem.} Test $Y\ci X\mid V$ with $V=(V_1,\dots,V_k)$ blocks.
\item \textbf{SO-PLS recap.} Sequential and orthogonalised PLS; order of blocks;
  number of components per block.
\item \textbf{Test statistic.} $T=\mathrm{RMSEP}(Y\sim V+X)-\mathrm{RMSEP}(Y\sim V)$ on
  held-out data (Monte Carlo cross-validation), $M_0$ and $M_1$ from one SO-PLS fit.
\item \textbf{Null distribution.} Resampling $X^*$ with the same MCCV splits:
  \begin{itemize}
  \item rotation of $X$ after linear adjustment for $V$ (Langs{\o}d-style rotation
    tests),
  \item global permutation,
  \item conditional permutation within groups of similar $V$.
  \end{itemize}
  Empirical and parametric (normal) p-value; ties and conservativeness when
  $b=0$ (\todo{randomised tie-breaking}).
\item \textbf{Choice of components.} Inner CV (auto) vs.\ one component per
  block vs.\ fixed; consequences for validity (Section 4).
\item \textbf{Alternative learners.} Multi-target XGBoost for nonlinear $V$ effects.
\item \textbf{Fast version.} Paired comparison of SO-PLS on real $X$ and on $X^*$
  over $K$ splits, resampled t-test; with and without Nadeau--Bengio correction.
\item \textbf{Residual-based tests.} Remove $V$ from $X$ and $Y$ with PLS or ROSA
  (cross-validated residuals), test dependence of residuals with RV, first canonical
  correlation or distance correlation, by (conditional) permutation.
\item \textbf{Theory.} Validity under linear Gaussian models; why it fails under
  nonlinear $V\to X, Y$; computational complexity.
  \todo{Proposition on exactness of the rotation test; discuss the residual-based
  tests under the same model.}
\item \textbf{Software.} \texttt{SOPLS.CCI.test}, \texttt{SOPLS.CCI.fast},
  \texttt{pls\_ci\_test}, \texttt{rosa\_ci\_test}.
\end{itemize}

% =============================================================================
\section{Simulation Study of the CI Tests}
\begin{itemize}
\item \textbf{Aim.} Type I error, power and run time of every test in Section 3.
\item \textbf{Design.}
  \begin{itemize}
  \item Scenarios (each with $Y\ci X\mid Z_1,Z_2$ true): common causes, chain,
    nearly collinear $X$, nonlinear/heteroskedastic/Laplace; plus the
    latent-proxy design as a case where the null is \emph{false}.
  \item Power: add a direct $X\to Y$ effect of increasing size, also a
    conditioning set that omits a confounder (e.g.\ $Y\sim X\mid Z_1$).
  \item Factors: $n\in\{200,500,1000,2000\}$, block dimensions, number of
    conditioning blocks, number of resamples $R$.
  \item Tests: SO-PLS MC (rotation/permutation/conditional), SO-PLS fast (with/without
    correction), XGBoost learner, PLS and ROSA residual tests (simple/conditional
    permutation, RV/cc1/dcor), and a baseline (e.g.\ Fisher-$z$ on canonical
    scores, \texttt{CCI}).
  \item Outcomes: rejection rate at $\alpha=0.05$ with Monte Carlo standard error,
    p-value uniformity under $H_0$ (QQ plots), power curves, run time.
  \end{itemize}
\item \textbf{Preliminary results} (linear $n=300$, 80 replications, $\alpha=0.05$):
  \begin{itemize}
  \item SO-PLS MC rotation: type I error 0.06; power 0.36/0.74 at effects 0.3/0.6.
  \item SO-PLS MC permutation: type I error 0.90 with one component per $V$ block
    (to be explained by the component rule); \todo{rerun with CV-chosen components}.
  \item Fast SO-PLS: 0.05 without correction (power 0.58/0.61), 0.00 with
    Nadeau--Bengio (power 0.23/0.48).
  \item PLS/ROSA residual tests: 0.05--0.08, power 1.00 at both effects (effects
    too large to separate the tests; \todo{use smaller effects}); 0.1--0.6 s per test.
  \item Nonlinear scenario: all tests reject 0.68--1.00 of true nulls.
  \end{itemize}
\item \textbf{To do.} Smaller effects, $n$ grid, larger $R$, more replications,
  multiple $V$ blocks, XGBoost and conditional resampling, nonlinear fix
  (polynomial expansion of $V$).
\end{itemize}

% =============================================================================
\section{Order of Blocks in SO-PLS}
\begin{itemize}
\item \textbf{Why order matters.} SO-PLS fits blocks sequentially and orthogonalises;
  the extra information of $X$ depends on what has been fitted before it.
\item \textbf{In the CI test.} $V$ before $X$ is required for the statistic. Order
  \emph{among} the blocks $V_1,\dots,V_k$ changes the component allocation and
  hence the fit, but not the hypothesis.
\item \textbf{Questions.}
  \begin{itemize}
  \item Does the order of $Z_1,Z_2$ change type I error or power?
  \item Strategies: fixed order, all orders and combine (min p / average),
    cbind blocks, order by explained variance.
  \item Interaction with the number of components per block.
  \end{itemize}
\item \textbf{Simulation} with permuted block orders on the scenarios of Section 4.
\item \textbf{Recommendation} for the default in the package.
\end{itemize}

% =============================================================================
\section{Causal Discovery Between Blocks}
\begin{itemize}
\item \textbf{PC algorithm at block level.}
  \begin{itemize}
  \item Skeleton: start from the complete graph over blocks, delete edge $(i,j)$ when
    $B_i\ci B_j\mid B_S$ for some $S$ of increasing size, store sepsets.
  \item Orient v-structures $B_i\to B_k\leftarrow B_j$ if $B_k\notin\mathrm{sepset}(i,j)$.
  \item Meek rules R1--R3 give a CPDAG over blocks.
  \end{itemize}
\item \textbf{Block-specific issues.}
  \begin{itemize}
  \item Conditioning sets are sets of blocks; number of tests is small (few
    nodes) but each test is costly $\Rightarrow$ choice between the fast and the
    Monte Carlo test.
  \item Different block sizes $p_k$ and different signal strengths $\Rightarrow$
    unequal power across edges.
  \item Effect of test errors (dense skeleton when the test over-rejects, missing
    edges with low power) on orientation.
  \item Stable / conservative PC variants; choice of $\alpha$; multiple testing.
  \end{itemize}
\item \textbf{Implementation.} Plug the block CI test into \texttt{pcalg::pc} through
  a user-supplied \texttt{indepTest}, or into the package's own PC code.
\item \textbf{Theory.} Soundness and completeness of PC with an oracle block CI test
  (follows from Section 2); consistency given a consistent test \todo{requirement
  on the test's size/power as $n\to\infty$}.
\item \textbf{Extensions.} FCI for hidden common causes between blocks; setting B
  via invariant causal prediction.
\end{itemize}

% =============================================================================
\section{Simulation}
\begin{itemize}
\item \textbf{Aim.} Quality of the recovered block graph.
\item \textbf{Design.}
  \begin{itemize}
  \item Random block DAGs ($K=4,6,8$ blocks, varying density and block sizes),
    data generated from latent factors per block and linear (then nonlinear)
    block-to-block effects.
  \item Factors: $n$, $\alpha$, signal strength, test used inside PC
    (SO-PLS MC, fast, PLS/ROSA residual).
  \item Competitors: variable-level PC on all columns, then collapse to blocks;
    PC on block summaries (first PC / canonical scores) with Fisher-$z$; the
    baseline methods in \texttt{R/discover.R}.
  \end{itemize}
\item \textbf{Metrics.} Skeleton precision and recall, orientation accuracy,
  structural Hamming distance (SHD) to the true CPDAG, run time.
\item \textbf{Baseline already available} (setting A, $n=1000$, 30 seeds, Fisher-$z$ on the
  baseline method): precision 0.84, recall 0.74, SHD 1.13.
\item \textbf{Robustness.} Nonlinearity, non-Gaussian noise, collinear blocks,
  violated sufficiency.
\end{itemize}

% =============================================================================
\section{Real Examples}
\begin{itemize}
\item \todo{Select two or three multiblock data sets where domain knowledge gives
  at least a partial ground truth (e.g.\ sensory/instrumental/consumer data,
  process-and-spectral data, omics layers).}
\item For each: data description, block definition, preprocessing, tests used,
  estimated block CPDAG, comparison with domain knowledge and with variable-level
  PC, sensitivity to $\alpha$ and to the test.
\item Discussion of what the discovered structure means and where it is implausible.
\end{itemize}

% =============================================================================
\section{Conclusions}
\begin{itemize}
\item Summary of contributions and main findings (validity and power of the test,
  usefulness for block-level discovery).
\item \textbf{Limitations.} Nonlinear dependence; faithfulness and sufficiency;
  tie-induced conservativeness; cost of the Monte Carlo test.
\item \textbf{Future work.} Nonlinear extension of the test, randomised p-values,
  FCI for blocks, setting B (environments), block-level effect estimation, further
  learners.
\end{itemize}

\appendix
\section*{Appendix}
\begin{itemize}
\item Proofs for Section 2 and for the rotation test.
\item Details of the simulators and additional tables and figures.
\item Software and reproducibility.
\end{itemize}

\begin{thebibliography}{99}

\bibitem{SG91} Spirtes, P. and Glymour, C. (1991). An algorithm for fast recovery of sparse causal graphs. \emph{Social Science Computer Review} 9(1), 62--72.
\bibitem{SGS00} Spirtes, P., Glymour, C. and Scheines, R. (2000). \emph{Causation, Prediction, and Search}, 2nd ed. MIT Press.
\bibitem{Pearl09} Pearl, J. (2009). \emph{Causality}, 2nd ed. Cambridge University Press.
\bibitem{Meek95} Meek, C. (1995). Causal inference and causal explanation with background knowledge. UAI.
\bibitem{KB07} Kalisch, M. and B\"uhlmann, P. (2007). Estimating high-dimensional directed acyclic graphs with the PC-algorithm. \emph{JMLR} 8, 613--636.
\bibitem{CM14} Colombo, D. and Maathuis, M. H. (2014). Order-independent constraint-based causal structure learning. \emph{JMLR} 15, 3741--3782.
\bibitem{SP20} Shah, R. D. and Peters, J. (2020). The hardness of conditional independence testing and the generalised covariance measure. \emph{Annals of Statistics} 48(3), 1514--1538.
\bibitem{Lang05} Langs{\o}d, \O. (2005). Rotation tests. \emph{Statistics and Computing} 15, 53--60.
\bibitem{NB03} Nadeau, C. and Bengio, Y. (2003). Inference for the generalization error. \emph{Machine Learning} 52, 239--281.
\bibitem{Naes11} N\ae s, T., Tomic, O., Mevik, B.-H. and Martens, H. (2011). Path modelling by sequential PLS regression. \emph{Journal of Chemometrics} 25, 28--40.
\bibitem{Lil16} Liland, K. H., N\ae s, T. and Indahl, U. G. (2016). ROSA --- a fast extension of partial least squares regression for multiblock data analysis. \emph{Journal of Chemometrics} 30, 651--662.
\bibitem{PBM16} Peters, J., B\"uhlmann, P. and Meinshausen, N. (2016). Causal inference by using invariant prediction. \emph{JRSS B} 78, 947--1012.
\bibitem{CCI} \todo{Thorjussen et al.: CCI test (thesis, Article 1) --- add full reference.}
\bibitem{SOPLS} \todo{J{\o}rgensen, Mevik and N\ae s (2007), SO-PLS --- check reference details.}

\end{thebibliography}
\end{document}
